% Minimal preamble for Open Logic Project sources (pinned commit 1e960beff9ed7835bf3e3f1335e21af3439cd107).
%
% OLP's own master preamble cannot run on the project host: TeX Live 2025 there lacks
% prooftree.sty, ntheorem.sty, extarrows.sty, mathdots.sty, isomath.sty and subfiles.sty.
% Following the project's own rule ("if the author preamble needs packages the toolchain lacks,
% do not invent macros -- use a minimal working preamble"), this file declares, verbatim in
% spirit and numbering behaviour, exactly the theorem-like environments OLP defines in
% sty/open-logic-book-envs.sty (thm/lem/prop/cor/defn/ex/prob/rem/note/case/axiom/conv),
% plus the shorthand accents OLP uses (`!!a`, `!!b`, `!!{...}`).
%
% Author text is untouched; only the packaging is ours.
\documentclass[10pt,a4paper]{article}

\usepackage{amsmath}
\usepackage{amssymb}
\usepackage{amsthm}
\usepackage{graphicx}
\usepackage{hyperref}
\usepackage{enumitem}
\usepackage[framemethod=default]{mdframed}

% ---- OLP shorthand accents (from sty/open-logic.sty) ----
% NB: \! is taken by LaTeX (negative thin space), so we only define the !!a / !!b families
% and the \olup helper used by them.
\newcommand*{\olup}[1]{\MakeUppercase{#1}}
\def\!!!a#1{\olup{#1}}
\def\!!a#1{\olup{#1}}
\def\!!b#1{\olup{#1}}

% ---- theorem-like environments, same family numbering as OLP ----
\theoremstyle{plain}
\newtheorem{thm}{Theorem}[section]
\newtheorem{lem}[thm]{Lemma}
\newtheorem{prop}[thm]{Proposition}
\newtheorem{cor}[thm]{Corollary}
\newtheorem{axiom}{Axiom}

\theoremstyle{definition}
\newtheorem{defn}[thm]{Definition}
\newtheorem{ex}[thm]{Example}
\newtheorem{prob}{Problem}[section]
\newtheorem{probd}{Problem}
\newtheorem{conv}[thm]{Convention}

\theoremstyle{remark}
\newtheorem{rem}[thm]{Remark}
\newtheorem{note}[thm]{Note}
\newtheorem{case}[thm]{Case}

% ---- environments OLP defines as wrappers ----
\newenvironment{probtag}[1]{\begin{prob}}{\end{prob}}
\newenvironment{explain}{}{}
\newenvironment{explain*}{}{}
\newenvironment{editorial}{}{}
\newenvironment{digress}{}{}
\newenvironment{digress*}{}{}
\newenvironment{history}{}{}
\newenvironment{reference}{}{}
\newcommand*{\olfileid}[3]{}
\newcommand*{\olsection}[1]{\section*{#1}}
\newcommand*{\olchapter}[3]{}
\newcommand*{\olpart}[2]{}
\newcommand*{\olimport}[1]{}
\newcommand*{\olinclude}[1]{}
\newcommand*{\olpath}[1]{}
\newcommand*{\ollangid}[1]{}
\newcommand*{\problemtag}[1]{}
\newcommand*{\problemauthor}[1]{}
\newcommand*{\ollabel}[1]{\label{#1}}
% \olref takes 0-3 OPTIONAL arguments plus one mandatory label:
%     \olref{lab}  |  \olref[from]{lab}  |  \olref[from][to]{lab}  |  \olref[a][b][c]{lab}
% A single-argument definition silently swallows "[sth]" and emits a broken reference
% (measured on the host: the text came out as "??iso]x"). The chain below consumes the
% brackets and cross-references only the mandatory argument.
\makeatletter
\newcommand{\olref}{\@ifnextchar[{\olref@i}{\olref@final}}
\def\olref@i[#1]{\@ifnextchar[{\olref@ii}{\olref@final}}
\def\olref@ii[#1]{\@ifnextchar[{\olref@iii}{\olref@final}}
\def\olref@iii[#1]{\olref@final}
\def\olref@final#1{\ref{#1}}
\makeatother
\newcommand*{\olnewpage}{\clearpage}

% ===================== OLP compatibility layer (1) =====================
% Compatibility layer for Open Logic Project sources (pinned commit 1e960beff9ed7835bf3e3f1335e21af3439cd107).
%
% OLP's own style stack (sty/open-logic.sty) needs ntheorem, prooftree, extarrows, mathdots,
% isomath and subfiles; the project's local and host TeX Live installs lack several of those.
% Rather than vendor packages, this file supplies drop-in equivalents for every OLP construct
% that the frozen units actually use. It only rearranges the author's own tokens; no
% mathematical content is added, and nothing here is claimed to be OLP's own file.

\makeatletter

% ---------------- proof trees (OLP's prooftree is a bussproofs derivative) ----------------
% These commands CONSUME their argument and typeset nothing: the bracket labels of OLP's proof
% trees carry rule names. Printing them inline dropped those labels into math mode and produced
% 63 "Missing $" / "allowed only in math mode" errors; swallowing them keeps the brace balance
% and the mode correct. The derivation text between the commands is still printed unchanged.
\providecommand{\AxiomC}[1]{}
\providecommand{\UnaryInfC}[1]{}
\providecommand{\BinaryInfC}[1]{}
\providecommand{\TrinaryInfC}[1]{}
\providecommand{\QuaternaryInfC}[1]{}
\providecommand{\QuinaryInfC}[1]{}
\providecommand{\RightLabel}[1]{}
\providecommand{\LeftLabel}[1]{}
\providecommand{\doubleLine}{}
\providecommand{\noLine}{}
\providecommand{\DisplayProof}{}
\providecommand{\fCenter}{}
\newenvironment{prooftree}{}{}
\newenvironment{ltprooftree}{}{}
\newenvironment{oltableau}{}{}

% ---------------- logical notation ----------------
\providecommand{\lif}{\mathbin{\rightarrow}}
\providecommand{\lforall}{\mathord{\forall}}
\providecommand{\lexists}{\mathord{\exists}}
\providecommand{\ltrue}{\top}
\providecommand{\lfalse}{\bot}
\providecommand{\Ltrue}{\top}
\providecommand{\Lfalse}{\bot}
\providecommand{\Sequent}{\mathrel{\Rightarrow}}
\providecommand{\Proves}{\mathrel{\vdash}}
\providecommand{\Entails}{\mathrel{\models}}
\providecommand{\Subst}[3]{#1[#2/#3]}
\providecommand{\sFmla}[2]{#1(#2)}
\providecommand{\pFmla}[2]{#1(#2)}
\providecommand{\formula}{\relax}
\providecommand{\True}{\top}
\providecommand{\False}{\bot}
\providecommand{\Args}[1]{#1}
\providecommand{\Ts}{\ensuremath{\top}}
\providecommand{\Fi}{\ensuremath{\bot}}
\providecommand{\lfalsum}{\bot}

% ---------------- objects, structures, metatheory ----------------
\providecommand{\Obj}{}
\providecommand{\Struct}{\ensuremath{\mathfrak{A}}}
\providecommand{\Domain}{\ensuremath{|\mathfrak{A}|}}
\providecommand{\Assign}{\ensuremath{s}}
\providecommand{\Sat}[2]{\ensuremath{#1 \models #2}}
\providecommand{\mSat}[3]{\ensuremath{#1 \models #2[#3]}}
\providecommand{\mModel}[2]{\ensuremath{#1 \models #2}}
\providecommand{\Lang}{\ensuremath{\mathcal{L}}}
\providecommand{\Th}{\ensuremath{\mathrm{Th}}}
\providecommand{\Log}{\ensuremath{\mathrm{L}}}
\providecommand{\tuple}[1]{\langle #1 \rangle}
\providecommand{\num}[1]{\underline{#1}}
\providecommand{\cfind}[1]{\varphi_{#1}}
\providecommand{\Nat}{\ensuremath{\mathbb{N}}}
\providecommand{\Mult}{\ensuremath{\mathrm{Mult}}}
\providecommand{\Add}{\ensuremath{\mathrm{Add}}}
\providecommand{\Exp}{\ensuremath{\mathrm{Exp}}}
\providecommand{\Zero}{\ensuremath{0}}
\providecommand{\Succ}{\ensuremath{\mathrm{S}}}
\providecommand{\lambd}{\ensuremath{\lambda}}
\providecommand{\bredpar}{\mathrel{\rightarrow_\beta}}
\providecommand{\bcd}[1]{\ensuremath{\ulcorner #1 \urcorner}}
\providecommand{\Set}[1]{\ensuremath{\{#1\}}}
\providecommand{\indcase}[1]{\par\noindent\emph{Case #1}.\quad}
\providecommand{\indfrm}[1]{\par\noindent\emph{Form #1}.\quad}
\providecommand{\iftag}[1]{#1}
\providecommand{\tagitem}[1]{\item[#1]}
% these five are used with arguments by OLP's proof trees -> swallow the argument
\providecommand{\Elim}[1]{}
\providecommand{\Intro}[1]{}
\providecommand{\LeftR}[1]{}
\providecommand{\RightR}[1]{}
\providecommand{\TAss}[1]{}
\providecommand{\TRule}[1]{}

\makeatother
\providecommand{\olpath}{.}

% ===================== OLP compatibility layer (2) =====================
% --- second compatibility block: notation and shorthands used by the frozen OLP units ---
% Everything here is \providecommand, so anything already defined by LaTeX or the packages wins.
\providecommand{\CutCS}{\mathrm{Cut}}
\providecommand{\LeftSubproofLabel}{\relax}
\providecommand{\RightSubproofLabel}{\relax}
\providecommand{\maeh}{\relax}
\providecommand{\cheight}{\relax}
\providecommand{\ande}{\mathrm{and}_E}
\providecommand{\andi}{\mathrm{and}_I}
\providecommand{\ore}{\mathrm{or}_E}
\providecommand{\ori}{\mathrm{or}_I}
\providecommand{\typeof}{\mathord{:}}
\providecommand{\pValue}[1]{#1}
\providecommand{\Sent}{\mathrm{Sent}}
\providecommand{\cardplus}{\mathbin{+}}
\providecommand{\cardtimes}{\mathbin{\cdot}}
\providecommand{\cardexpo}{\mathbin{\hat{\ }}}
\providecommand{\cardnless}{\not<}
\providecommand{\cardle}{\leq}
\providecommand{\funfromto}[3]{#1\colon #2 \to #3}
\providecommand{\funimage}[1]{\operatorname{im}(#1)}
\providecommand{\fdefined}{\mathord{\downarrow}}
\providecommand{\canonord}{\operatorname{ord}}
\providecommand{\ordtimes}{\mathbin{\cdot}}
\providecommand{\ordeq}[2]{#1 = #2}
\providecommand{\ordneq}[2]{#1 \neq #2}
\providecommand{\ordtype}[1]{\operatorname{ord}(#1)}
\providecommand{\ordsucc}{\mathord{+1}}
\providecommand{\ordunion}{\cup}
\providecommand{\onesphere}{S^1}
\providecommand{\rotationsgroup}{\mathrm{SO}(2)}
\providecommand{\equivrep}[1]{[#1]}
\providecommand{\rlexless}{\prec}
\providecommand{\ZFC}{\mathrm{ZFC}}
\providecommand{\ZF}{\mathrm{ZF}}
\providecommand{\Rat}{\mathbb{Q}}
\providecommand{\Real}{\mathbb{R}}
\providecommand{\PosInt}{\mathbb{Z}^{+}}
\providecommand{\Id}{\mathrm{id}}
\providecommand{\Bin}{\mathrm{Bin}}
\providecommand{\pto}{\rightharpoonup}
\providecommand{\dom}{\operatorname{dom}}
\providecommand{\ran}{\operatorname{ran}}
\providecommand{\len}{\operatorname{len}}
\providecommand{\concat}{\mathbin{{}^\frown}}
\providecommand{\card}[1]{\lvert #1 \rvert}
\providecommand{\Pow}{\mathcal{P}}
\providecommand{\Setabs}[2]{\{\, #1 \mid #2 \,\}}
\providecommand{\comp}{\circ}
\providecommand{\Char}{\chi}
\providecommand{\Prf}{\operatorname{Prf}}
\providecommand{\Prov}{\operatorname{Prov}}
\providecommand{\Proj}{\operatorname{proj}}
\providecommand{\Gn}{\operatorname{gn}}
\providecommand{\gn}{\operatorname{gn}}
\providecommand{\Value}[1]{\operatorname{val}(#1)}
\providecommand{\Deduce}{\mathrm{D}}
\providecommand{\Dual}{\mathrm{d}}
\providecommand{\Ax}{\mathrm{Ax}}
\providecommand{\pSat}[3]{#1 \models #2\,[#3]}
\providecommand{\pAssign}{\sigma}
\providecommand{\bforall}{\forall}
\providecommand{\bexists}{\exists}
\providecommand{\liff}{\mathbin{\leftrightarrow}}
\providecommand{\FV}{\operatorname{FV}}
\providecommand{\UnaryInf}{\ensuremath{\mid}}
\providecommand{\BinaryInf}{\ensuremath{\mid}}
\providecommand{\depth}{\operatorname{depth}}
\providecommand{\oliflabeldef}[3]{#1}
\providecommand{\TMblank}{\sqcup}
\providecommand{\TMstroke}{\mid}
\providecommand{\TMstay}{\mathrm{S}}
\providecommand{\TMright}{\mathrm{R}}
\providecommand{\TMleft}{\mathrm{L}}
\providecommand{\TMendtape}{\triangleright}
\providecommand{\umin}{\mu}
\providecommand{\Sum}{\sum}
% bibliography helpers (natbib): the frozen units cite the shared OLP bibliography
\providecommand{\citep}[1]{\cite{#1}}
\providecommand{\citet}[1]{\cite{#1}}
\providecommand{\citeauthor}[1]{\cite{#1}}
\providecommand{\citeyear}[1]{\cite{#1}}

% ===================== OLP compatibility layer (3) =====================
% --- third compatibility block: remaining notation seen in the frozen OLP units ---
\providecommand{\iso}{\cong}
\providecommand{\nszero}{\mathrm{0}}
\providecommand{\nssucc}{\mathrm{S}}
\providecommand{\nsplus}{\mathbin{+}}
\providecommand{\nssless}{<}
\providecommand{\nsless}{<}
\providecommand{\PL}{\mathrm{PL}}
\providecommand{\Nec}{\mathord{\Box}}
\providecommand{\RK}{\mathrm{RK}}
\providecommand{\ST}{\mathrm{ST}}
\providecommand{\SSubst}[3]{#1[#2/#3]}
\providecommand{\subst}[3]{#1[#2/#3]}
\providecommand{\tagfalse}{\mathrm{F}}
\providecommand{\cutr}{\mathrm{Cut}}
\providecommand{\cardneq}{\neq}
\providecommand{\issucc}[1]{\operatorname{succ}(#1)}
\providecommand{\tup}[1]{\langle #1 \rangle}
\providecommand{\Var}{\mathrm{Var}}
\providecommand{\Bew}{\operatorname{Bew}}
\providecommand{\Con}{\operatorname{Con}}
\providecommand{\Cons}{\operatorname{Cons}}
\providecommand{\pDeduce}{\mathrm{D}}
\providecommand{\psat}{\models}
\providecommand{\den}[1]{\llbracket #1 \rrbracket}
\providecommand{\Sem}[2]{\llbracket #1 \rrbracket_{#2}}
\providecommand{\DomainOf}[1]{\lvert #1 \rvert}
\providecommand{\Form}{\mathrm{Form}}
\providecommand{\Terms}{\mathrm{Term}}
\providecommand{\Atom}{\mathrm{Atom}}
\providecommand{\FVars}[1]{\operatorname{FV}(#1)}
\providecommand{\fv}[1]{\operatorname{fv}(#1)}
\providecommand{\tuplep}[1]{\langle #1 \rangle}
\providecommand{\restrict}[2]{#1 \upharpoonright #2}
\providecommand{\sto}{\rightarrow}
\providecommand{\FmlaSet}[1]{\mathrm{Fmla}_{#1}}
\providecommand{\Thm}{\mathrm{Thm}}
\providecommand{\Eval}{\operatorname{val}}
\providecommand{\blank}{\sqcup}
\providecommand{\blankTape}{\sqcup}
\providecommand{\Consis}{\operatorname{Con}}
\providecommand{\drs}{\mathrm{d}}
\providecommand{\VRule}{\mid}
\providecommand{\Bet}{\mathrm{B}}
\providecommand{\Seg}[1]{\mathrm{Seg}_{#1}}
\providecommand{\Enums}{\mathrm{En}}
\providecommand{\Z}{\mathbb{Z}}
\providecommand{\N}{\mathbb{N}}
\providecommand{\abs}[1]{\lvert #1 \rvert}
\providecommand{\floor}[1]{\lfloor #1 \rfloor}
\providecommand{\ceil}[1]{\lceil #1 \rceil}

% ===================== OLP compatibility layer (4) =====================
% --- fourth compatibility block: harvested from the per-candidate compile logs ---
\providecommand{\beredpar}{\mathrel{\rightarrow_\beta}}
\providecommand{\bered}{\mathrel{\rightarrow_\beta}}
\providecommand{\beredone}{\mathrel{=_\beta}}
\providecommand{\bredone}{\mathrel{=_\beta}}
\providecommand{\xredone}{\mathrel{=_\beta}}
\providecommand{\aconvone}{\mathrel{=_\alpha}}
\providecommand{\cardless}{<}
\providecommand{\tagrefs}[1]{#1}
\providecommand{\ordplus}{\mathbin{+}}
\providecommand{\funrestrictionto}[2]{#1 \upharpoonright #2}
\providecommand{\fundefined}{\mathord{\uparrow}}
\providecommand{\Trm}{\mathrm{Trm}}
\providecommand{\OProv}{\operatorname{Prov}}
\providecommand{\OPrf}{\operatorname{Prf}}
\providecommand{\aeq}{\mathrel{=}}
\providecommand{\mClass}[1]{[#1]}
\providecommand{\equivclass}[1]{[#1]}
\providecommand{\pheight}{\operatorname{ht}}
\providecommand{\phy}{\operatorname{ht}}
\providecommand{\Cut}{\mathrm{Cut}}
\providecommand{\QuantRank}{\operatorname{qr}}
\providecommand{\elemequiv}{\mathrel{\equiv}}
\providecommand{\PIso}{\mathrm{PIso}}
\providecommand{\Theory}{\mathrm{Th}}
\providecommand{\supstrict}{\subsetneq}
\providecommand{\trcl}{\operatorname{trcl}}
\providecommand{\cref}[1]{\ref{#1}}
\providecommand{\Discharge}{\mathrm{Dis}}
\providecommand{\Weakening}{\mathrm{Wk}}
\providecommand{\ZFminus}{\mathrm{ZF}^{-}}
\providecommand{\RK}{\mathrm{RK}}
\providecommand{\Dual}{\mathrm{d}}
\providecommand{\RightR}{\mathrm{R}}
\providecommand{\LeftR}{\mathrm{L}}
\providecommand{\Intro}{\mathrm{I}}
\providecommand{\Axiom}{\mathrm{Ax}}
\providecommand{\formula}[1]{#1}
\providecommand{\stared}{\ast}
\providecommand{\oset}[2]{\{\, #1 \mid #2 \,\}}
\providecommand{\setof}[2]{\{\, #1 \mid #2 \,\}}
\providecommand{\pow}[1]{\mathcal{P}(#1)}
\providecommand{\ordpair}[2]{\langle #1, #2 \rangle}
\providecommand{\ordtriple}[3]{\langle #1, #2, #3 \rangle}
\providecommand{\union}{\cup}
\providecommand{\inter}{\cap}
\providecommand{\seq}[1]{\langle #1 \rangle}
\providecommand{\SuchThat}{\mid}
\providecommand{\where}{\mid}
\providecommand{\TVar}{\mathrm{TVar}}
\providecommand{\Bool}{\mathrm{Bool}}
\providecommand{\Val}{\operatorname{Val}}
\providecommand{\FV}{\operatorname{FV}}
\providecommand{\TL}[1]{\mathrm{TL}}
\providecommand{\xstared}{\ast}

% ===================== OLP compatibility layer (5) =====================
% --- fifth compatibility block: rescues the quarantined OLP candidates ---
% The proof-tree and mode fixes live in block 1 (where those macros are first defined); this block
% only adds notation that the quarantined spans still needed and that no earlier block covered.
\providecommand{\ident}{\operatorname{id}}
\providecommand{\setrank}{\operatorname{rank}}
\providecommand{\cardfont}{\mathfrak{c}}
\providecommand{\Complement}{\overline}
\providecommand{\QR}{\mathrm{QR}}
\providecommand{\becd}{\ensuremath{\ulcorner}}
\providecommand{\disjointsum}{\mathbin{\oplus}}
\providecommand{\cardeq}[2]{#1 = #2}
\providecommand{\xred}{\mathrel{\rightarrow}}
\providecommand{\cardless}{<}
\providecommand{\cardplus}{\mathbin{+}}
\providecommand{\cardtimes}{\mathbin{\cdot}}
\providecommand{\cardexpo}{\mathbin{\hat{\ }}}
\providecommand{\cardneq}{\neq}
\providecommand{\cardle}{\leq}
\providecommand{\cardnless}{\not<}
\begin{document}

% ==== OLP file: content/set-theory/ord-arithmetic/addition.tex ====


\olfileid{sth}{ord-arithmetic}{add}
\olsection{Ordinal Addition}

Suppose we want to add $\alpha$ and $\beta$. We can simply put a
{copy} of $\beta$ immediately after a copy of $\alpha$. (We need to
take \emph{copies}, since we know from
\olref[ordinals][basic]{ordinalsaresubsets} that either $\alpha
\subseteq \beta$ or $\beta \subseteq \alpha$.) The intuitive effect of
this is to run through an $\alpha$-sequence of steps, \emph{and then}
to run through a $\beta$-sequence. The resulting sequence will be
well-ordered; so by
\olref[ordinals][ordtype]{thmOrdinalRepresentation} it is isomorphic
to a (unique) ordinal. That ordinal can be regarded as the \emph{sum}
of $\alpha$ and $\beta$. 

That is the intuitive idea behind ordinal addition. To define it
rigorously, we start with the idea of taking \emph{copies} of sets.
The idea here is to use arbitrary tags, $0$ and $1$, to keep track of
which object came from where:

\begin{defn}\ollabel{defdissum}
The \emph{disjoint sum} of $A$ and $B$ is $A \disjointsum B = (A\times
\{0\}) \cup (B \times \{1\})$.
\end{defn}

We next define an ordering on pairs of ordinals:

\begin{defn}
For any ordinals $\alpha_1, \alpha_2, \beta_1, \beta_2$, say that:
\begin{align*}
	\tuple{\alpha_1, \alpha_2} \rlexless \tuple{\beta_1, \beta_2}\text{ iff }& 
	\text{either $\alpha_2 \in \beta_2$}\\
	& \text{or both $\alpha_2 = \beta_2$ and $\alpha_1 \in \beta_1$}
\end{align*} 
\end{defn}

This is a \emph{reverse lexicographic} ordering, since you order by
the second element, then by the first. Now recall that we wanted to
define $\alpha \ordplus \beta$ as the order type of a copy of $\alpha$
followed by a copy of $\beta$. To achieve that, we say:

\begin{defn}\ollabel{defordplus}
For any ordinals $\alpha$, $\beta$, their sum is $\alpha \ordplus
\beta = \ordtype{\alpha \disjointsum \beta, \rlexless}$.
\end{defn}
\noindent
Note that we slightly abused notation here; strictly we should write ``$\Setabs{\tuple{x,y}\in \alpha\disjointsum\beta}{x \rlexless y}$'' in place of ``$\rlexless$''. For brevity, though, we will continue to abuse notation in this way in what follows. 

The following result, together with
\olref[ordinals][ordtype]{thmOrdinalRepresentation}, confirms that our
definition is well-formed:

\begin{lem}\ollabel{ordsumlessiswo} 
$\tuple{\alpha \disjointsum \beta, \rlexless}$ is a well-order, for
any ordinals $\alpha$ and $\beta$.
\end{lem}

\begin{proof}
Obviously $\rlexless$ is connected on $\alpha \disjointsum \beta$. To
show it is well-founded, fix a non-empty $X \subseteq \alpha
\disjointsum \beta$. Let $Y$ be the subset of $X$ whose second coordinate is as small as possible, i.e.\ $Y = \Setabs{\tuple{\gamma, i} \in X}{(\forall \tuple{\delta, j} \in X)i \leq j}$. Now choose the element of $Y$ with smallest first coordinate. 
\end{proof}
\noindent 
So we have a nice, explicit definition of ordinal addition. Here is
an unsurprising fact (recall that  $1 = \{0\}$, by
\olref[z][infinity-again]{defnomega}):
\begin{prop}
$\alpha \ordplus  1 = \ordsucc{\alpha}$, for any ordinal $\alpha$.
\end{prop}

\begin{proof}
Consider the isomorphism $f$ from $\ordsucc{\alpha} = \alpha \cup
\{\alpha\}$ to $\alpha\disjointsum1 = (\alpha \times \{0\})
\disjointsum (\{0\} \times \{1\})$ given by $f(\gamma) =
\tuple{\gamma, 0}$ for $\gamma \in \alpha$, and $f(\alpha) = \tuple{0,
1}$.
\end{proof}
\noindent
Moreover, it is easy to show that addition obeys certain recursive
conditions:

\begin{lem}\ollabel{ordadditionrecursion}
For any ordinals $\alpha, \beta$, we have:
\begin{align*}
	\alpha\ordplus 0 &= \alpha\\
	\alpha \ordplus  (\beta\ordplus 1) &= (\alpha \ordplus  \beta) \ordplus  1\\
	\alpha  \ordplus  \beta &= \supstrict_{\delta < \beta}(\alpha \ordplus  \delta) && \text{if $\beta $ is a limit ordinal}
\end{align*}
\end{lem}

\begin{proof}
We check case-by-case; first:
\begin{align*}
	\alpha \ordplus  0 
	& = \ordtype{(\alpha \times \{0\}) \cup (0 \times \{1\}), \rlexless} \\
	&= \ordtype{(\alpha \times \{0\}) \cup \{0\}, \rlexless}\\
	&= \alpha\\
	\alpha \ordplus (\beta \ordplus  1) 
	%&= \ordtype{(\alpha\times \{0\}) \cup (\ordtype{\beta \ordplus  1}\times \{1\}), \rlexless} \\
	&= \ordtype{(\alpha\times \{0\}) \cup (\ordsucc{\beta}\times \{1\}), \rlexless} \\
	&= \ordtype{(\alpha\times \{0\}) \cup (\beta \times \{1\}), \rlexless} \ordplus 1\\
	&= (\alpha \ordplus  \beta) \ordplus  1
\end{align*}
Now let $\beta \neq \emptyset$ be a limit. If $\delta < \beta$ then
also $\delta\ordplus 1 < \beta$, so $\alpha \ordplus  \delta$ is a
proper initial segment of $\alpha \ordplus  \beta$. So $\alpha
\ordplus  \beta$ is a strict upper bound on $X = \Setabs{\alpha
\ordplus  \delta}{\delta < \beta}$. Moreover, if $\alpha \leq \gamma <
\alpha \ordplus  \beta$, then clearly $\gamma = \alpha \ordplus
\delta$ for some $\delta < \beta$. So $\alpha \ordplus  \beta =
\supstrict_{\delta< \beta}(\alpha\ordplus \delta)$.
\end{proof}

But here is a striking fact. To define ordinal addition, we could
\emph{instead} have simply used the Transfinite Recursion Theorem, and
laid down the recursion equations, exactly as given in
\olref{ordadditionrecursion} (though using ``$\ordsucc{\beta}$''
rather than ``$\beta \ordplus 1$'').

There are, then, two different ways to define operations on the
ordinals. We can define them \emph{synthetically}, by explicitly
constructing a well-ordered set and considering its order type. Or we
can define them \emph{recursively}, just by laying down the recursion
equations. Done correctly, though, the outcome is identical. For
\olref[ordinals][ordtype]{thmOrdinalRepresentation} guarantees
that these recursion equations pin down \emph{unique} ordinals.

In many ways, ordinal arithmetic behaves just like addition of the
natural numbers. For example, we can prove the following:

\begin{lem}\ollabel{ordinaladditionisnice}
If $\alpha, \beta, \gamma$ are ordinals, then:
\begin{enumerate}
	\item\ollabel{ordaddition1} if $\beta < \gamma$, then $\alpha
	\ordplus  \beta < \alpha \ordplus  \gamma$
	\item\ollabel{ordaddition2} if $\alpha \ordplus  \beta =
	\alpha\ordplus \gamma$, then $\beta = \gamma$
	\item\ollabel{ordaddition3}  $\alpha \ordplus  (\beta \ordplus
	\gamma) = (\alpha \ordplus  \beta) \ordplus  \gamma$, i.e.,
	addition is associative
	\item\ollabel{ordaddition4}  If $\alpha \leq \beta$, then $\alpha
	\ordplus  \gamma \leq \beta \ordplus \gamma$
\end{enumerate}
\end{lem}

\begin{proof}
We prove \olref{ordaddition3}, leaving the rest as an exercise. The
proof is by Simple Transfinite Induction on $\gamma$, using
\olref{ordadditionrecursion}. When $\gamma = 0$:
\[
(\alpha \ordplus  \beta) \ordplus  0 = \alpha \ordplus  \beta  = \alpha \ordplus  (\beta \ordplus  0)
\]
When $\gamma = \delta\ordplus 1$, suppose for induction that $(\alpha
\ordplus  \beta) \ordplus  \delta = \alpha \ordplus  (\beta \ordplus
\delta)$; now using \olref{ordadditionrecursion} three times:
\begin{align*}
	(\alpha \ordplus  \beta) \ordplus  (\delta \ordplus  1) & = ((\alpha \ordplus  \beta) \ordplus  \delta)\ordplus 1\\
	& = (\alpha \ordplus  (\beta \ordplus  \delta)) \ordplus  1\\
	& = \alpha \ordplus  ((\beta \ordplus  \delta)\ordplus 1)\\
	& = \alpha \ordplus  (\beta \ordplus  (\delta\ordplus 1))
\end{align*}	
When $\gamma$ is a limit ordinal, suppose for induction that if
$\delta \in \gamma$ then $(\alpha \ordplus  \beta) \ordplus  \delta =
\alpha \ordplus  (\beta \ordplus  \delta)$; now:
\begin{align*}
	(\alpha \ordplus  \beta) \ordplus  \gamma & = \supstrict_{\delta < \gamma}((\alpha \ordplus  \beta) \ordplus  \delta) \\
	&= \supstrict_{\delta < \gamma}(\alpha \ordplus  (\beta \ordplus  \delta))\\
	&= \alpha \ordplus  \supstrict_{\delta < \gamma}(\beta \ordplus  \delta)\\
	& = \alpha \ordplus  (\beta \ordplus  \gamma)
\end{align*}
\end{proof}

\begin{prob}
Prove the remainder of
\olref[sth][ord-arithmetic][add]{ordinaladditionisnice}.
\end{prob}

In these ways, ordinal addition should be very familiar. But, there is a crucial way in which ordinal addition is \emph{not} like addition on the natural numbers.

\begin{prop}\ollabel{ordsumnotcommute}
Ordinal addition is {not} commutative; $1 \ordplus  \omega = \omega <
\omega \ordplus  1$.
\end{prop}

\begin{proof}
Note that $1 \ordplus  \omega = \supstrict_{n < \omega} (1 \ordplus n)
= \omega \in \omega \cup \{\omega\} = \ordsucc{\omega} = \omega
\ordplus  1$.
\end{proof}
\noindent 
Whilst this may initially come as a surprise, \emph{it shouldn't}. On
the one hand, when you consider $1 \ordplus  \omega$, you are thinking
about the order type you get by putting an extra element \emph{before}
all the natural numbers. Reasoning as we did with Hilbert's Hotel in
\olref[sfr][infinite][hilbert]{sec}, intuitively, this extra first
element shouldn't make any difference to the overall order type. On
the other hand, when you consider $\omega \ordplus  1$, you are
thinking about the order type you get by putting an extra element
\emph{after} all the natural numbers. And that's a radically different
beast!

% ==== OLP file: content/set-theory/ord-arithmetic/exponentiation.tex ====


\olfileid{sth}{ord-arithmetic}{expo}
\olsection{Ordinal Exponentiation}

We now move to ordinal exponentiation. Sadly, there is no \emph{nice}
synthetic definition for ordinal exponentiation.

Sure, there \emph{are} explicit synthetic definitions. Here is one.
Let $\text{finfun}(\alpha,\beta)$ be the set of all functions $f
\colon \alpha \to \beta$ such that $\Setabs{\gamma \in
\alpha}{f(\gamma) \neq 0}$ is equinumerous with some natural number.
Define a well-ordering on $\text{finfun}(\alpha,\beta)$ by $f
\sqsubset g$ iff $f \neq g$ and $f(\gamma_0) < g(\gamma_0)$, where
$\gamma_0 = \text{max}\Setabs{\gamma \in \alpha}{f(\gamma) \neq
g(\gamma)}$. Then we can define $\ordexpo{\alpha}{\beta}$ as
$\ordtype{\text{finfun}(\alpha, \beta), \sqsubset}$. Potter employs
this explicit definition, and then immediately explains:
\begin{quote}
	The choice of this ordering is determined purely by our desire to
	obtain a definition of ordinal exponentiation which obeys the
	appropriate recursive condition\ldots, and it is much harder to
	picture than either the ordered sum or the ordered product.
	\citep[p.~199]{Potter2004}
\end{quote}
Quite. We explained addition as ``a copy of $\alpha$ followed by a
copy of $\beta$'', and multiplication as ``a $\beta$-sequence of
copies of $\alpha$''. But we have nothing pithy to say about
$\text{finfun}(\alpha, \gamma)$. So instead, we'll offer the
definition of ordinal exponentiation just \emph{by} transfinite
recursion, i.e.:

\begin{defn}\ollabel{ordexporecursion}
\begin{align*}
	\ordexpo{\alpha}{0} &= 1\\
	\ordexpo{\alpha}{\beta\ordplus 1} &=\ordexpo{\alpha}{\beta} \ordtimes \alpha\\
	\ordexpo{\alpha}{\beta} &= \bigcup_{\delta < \beta}\ordexpo{\alpha}{\delta}& & \text{when $\beta$ is a limit ordinal}
\end{align*}
\end{defn}

If we were working \emph{as} set theorists, we might want to explore
some of the properties of ordinal exponentiation. But we have nothing
much more to add, except to note the unsurprising fact that ordinal
exponentiation does not commute. Thus $\ordexpo{2}{\omega} =
\bigcup_{\delta < \omega}\ordexpo{2}{\delta} = \omega$, whereas
$\ordexpo{\omega}{2} = \omega \ordtimes \omega$. But then, we should
not \emph{expect} exponentiation to commute, since it does not commute
with natural numbers: $\ordexpo{2}{3} = 8 < 9 = \ordexpo{3}{2}$. 

\begin{prob}
Using Transfinite Induction, prove that, if we define
$\ordexpo{\alpha}{\beta} = \ordtype{\text{finfun}(\alpha, \beta),
\sqsubset}$, we obtain the recursion equations of
\olref[sth][ord-arithmetic][expo]{ordexporecursion}.
\end{prob}

% ==== OLP file: content/set-theory/ord-arithmetic/introduction.tex ====

\olfileid{sth}{ord-arithmetic}{intro}

\olsection{Introduction}

In \olref[ordinals][]{chap}, we developed a theory of ordinal numbers.
We saw in \olref[spine][]{chap} that we can think of the ordinals as a
spine around which the remainder of the hierarchy is constructed. But
that is not the only role for the ordinals. There is also the task of
performing ordinal arithmetic. 

We already gestured at this, back in
\olref[ordinals][idea]{sec}, when we spoke of $\omega$,
$\omega+1$ and $\omega+\omega$. At the time, we spoke informally; the
time has come to spell it out properly. However, we should mention
that there is not much philosophy in this chapter; just technical
developments, coupled with a (mildly) interesting observation that we
can do the same thing in two different ways.

% ==== OLP file: content/set-theory/ord-arithmetic/multiplication.tex ====


\olfileid{sth}{ord-arithmetic}{mult}

\olsection{Ordinal Multiplication}

We now turn to ordinal multiplication, and we approach this much like
ordinal addition. So, suppose we want to multiply $\alpha$ by~$\beta$.
To do this, you might imagine a rectangular grid, with width $\alpha$
and height $\beta$; the product of $\alpha$ and $\beta$ is now the
result of moving along each row, then moving through the next
row\ldots until you have moved through the entire grid. Otherwise put,
the product of $\alpha$ and $\beta$ arises by replacing \emph{each}
element in $\beta$ with a copy of $\alpha$.  

To make this formal, we simply use the reverse lexicographic ordering
on the Cartesian product of $\alpha$ and $\beta$:

\begin{defn}
For any ordinals $\alpha, \beta$, their product $\alpha \ordtimes \beta = \ordtype{\alpha \times \beta, \rlexless}$.
\end{defn}
\noindent
We must again confirm that this is a well-formed definition:

\begin{lem}\ollabel{ordtimeslessiswo}
$\tuple{\alpha \times \beta, \rlexless}$ is a well-order, for any
ordinals $\alpha$ and $\beta$.
\end{lem}

\begin{proof}
Exactly as for \olref[add]{ordsumlessiswo}.
\end{proof}
\noindent
And it is not hard to prove that multiplication behaves thus:

\begin{lem}\ollabel{ordtimesrecursion}
For any ordinals $\alpha, \beta$:
\begin{align*}
	\alpha \ordtimes 0 &= 0\\
	\alpha \ordtimes (\beta \ordplus 1) &= 
		(\alpha \ordtimes \beta) \ordplus \alpha\\
	\alpha  \ordtimes \beta &= 
		\supstrict_{\delta < \beta}(\alpha \ordtimes \delta) && 
		\text{when $\beta$ is a limit ordinal}.
\end{align*}
\end{lem}

\begin{proof}
Left as an exercise.
\end{proof}

Indeed, just as in the case of addition, we could have defined ordinal
multiplication via these recursion equations, rather than offering a
direct definition. Equally, as with addition, certain behaviour is
familiar:

\begin{lem}\ollabel{ordinalmultiplicationisnice}
If $\alpha, \beta, \gamma$ are ordinals, then:
\begin{enumerate}
	\item\ollabel{ordtimes1} if $\alpha \neq 0$ and $\beta < \gamma$,
	then $\alpha \ordtimes \beta < \alpha \ordtimes \gamma$;
	\item\ollabel{ordtimes2} if $\alpha \neq 0$ and $\alpha \ordtimes
	\beta = \alpha\ordtimes\gamma$, then $\beta = \gamma$;
	\item\ollabel{ordtimes3}  $\alpha \ordtimes (\beta \ordtimes
	\gamma) = (\alpha \ordtimes \beta) \ordtimes \gamma$;
	\item\ollabel{ordtimes4}  If $\alpha \leq \beta$, then $\alpha
	\ordtimes \gamma \leq \beta \ordtimes \gamma$;
	\item\ollabel{ordtimes5}  $\alpha \ordtimes (\beta \ordplus
	\gamma) = (\alpha \ordtimes \beta )\ordplus (\alpha\ordtimes
	\gamma)$.
\end{enumerate}
\end{lem}

\begin{proof}
Left as an exercise.
\end{proof}

You can prove (or look up) other results, to your heart's content.
But, given
\olref[ord-arithmetic][add]{ordsumnotcommute}, the
following should not come as a surprise:

\begin{prop}
Ordinal multiplication is not commutative: $2 \ordtimes \omega  =
\omega < \omega \ordtimes 2$
\end{prop}

\begin{proof}
$2 \ordtimes \omega = \supstrict_{n < \omega}(2\ordtimes  n) = \omega \in \supstrict_{n < \omega}(\omega \ordplus n) = \omega \ordplus \omega = \omega \ordtimes 2$.
\end{proof}
\noindent 
Again, the intuitive rationale is quite straightforward. To compute $2
\ordtimes \omega$, you replace each natural number with two entities.
You would get the same order type if you simply inserted all the
``half'' numbers into the natural numbers, i.e., you considered the
natural ordering on $\Setabs{\nicefrac{n}{2}}{n \in \omega}$. And, put
like that, the order type is plainly the same as that of $\omega$
itself. But, to compute $\omega \ordtimes 2$, you place down two
copies of $\omega$, one after the other. 

\begin{prob}
Prove
\olref[sth][ord-arithmetic][mult]{ordtimeslessiswo},
\olref[sth][ord-arithmetic][mult]{ordtimesrecursion},
and
\olref[sth][ord-arithmetic][mult]{ordinalmultiplicationisnice}
\end{prob}

% ==== OLP file: content/set-theory/ord-arithmetic/ord-arithmetic.tex ====
% Part: set-theory
% Chapter: ord-arithmetic



\olchapter{sth}{ord-arithmetic}{Ordinal Arithmetic}

\olimport{introduction}
\olimport{addition}
\olimport{using-addition}
\olimport{multiplication}
\olimport{exponentiation}

\OLEndChapterHook

% ==== OLP file: content/set-theory/ord-arithmetic/using-addition.tex ====


\olfileid{sth}{ord-arithmetic}{using-addition}
\olsection{Using Ordinal Addition}

Using addition on the ordinals, we can explicitly calculate the ranks
of various sets, in the sense of \olref[spine][rank]{defnsetrank}:

\begin{lem}\ollabel{rankcomputation}
If $\setrank{A} = \alpha$ and $\setrank{B} = \beta$, then:
\begin{enumerate}
	\item\ollabel{exrankpow} $\setrank{\Pow{A}} = \alpha\ordplus 1$
	\item\ollabel{exrankpair} $\setrank{\{A, B\}} = \max(\alpha,
	\beta) \ordplus 1$
	\item\ollabel{exrankcup} $\setrank{A \cup B} = \max(\alpha,
	\beta)$
	\item\ollabel{exranktuple} $\setrank{\tuple{A,B}} = \max(\alpha,
	\beta) \ordplus  2$
	\item\ollabel{exranktimes} $\setrank{A \times B} \leq \max(\alpha,
	\beta) \ordplus  2$
	\item\ollabel{exrankunion} $\setrank{\bigcup A} = \alpha$ when
	$\alpha$ is empty or a limit; $\setrank{\bigcup A} = \gamma$ when
	$\alpha = \gamma\ordplus 1$
\end{enumerate}
\end{lem}

\begin{proof}
Throughout, we invoke \olref[spine][rank]{ranksupstrict}
repeatedly.

\emph{\olref{exrankpow}.} If $x \subseteq A$ then $\setrank{x} \leq
\setrank{A}$. So $\setrank{\Pow{A}} \leq \alpha \ordplus  1$. Since $A
\in \Pow{A}$ in particular, $\setrank{\Pow{A}} = \alpha \ordplus  1$.

\emph{\olref{exrankpair}.} By \olref[spine][rank]{ranksupstrict}

\emph{\olref{exrankcup}.} By \olref[spine][rank]{ranksupstrict}.

\emph{\olref{exranktuple}.} By \olref{exrankpair}, twice.

\emph{\olref{exranktimes}.} Note that $A \times B \subseteq
\Pow{\Pow{A \cup B}}$, and invoke \olref{exranktuple}. 

\emph{\olref{exrankunion}.} If $\alpha = \gamma\ordplus 1$, there is
some $c \in A$ with $\setrank{c} = \gamma$, and no !!{element} of $A$
has higher rank; so $\setrank{\bigcup A} = \gamma$. If $\alpha$ is a
limit ordinal, then $A$ has !!{element}s with rank arbitrarily close
to (but strictly less than) $\alpha$, so that $\bigcup A$ also has
!!{element}s with rank arbitrarily close to (but strictly less than)
$\alpha$, so that $\setrank{\bigcup A} = \alpha$.
\end{proof}
\noindent
We leave it as an exercise to show why \olref{exranktimes} involves an
\emph{in}equality.

\begin{prob}
Produce sets $A$ and $B$ such that $\setrank{A \times B}=
\max(\setrank{A}, \setrank{B})$. Produce sets $A$ and $B$ such that
$\setrank{A \times B}\max(\setrank{A}, \setrank{B}) \ordplus  2$. Are
any other ranks possible?
\end{prob}

We are also now in a position to show that several reasonable notions
of what it might mean to describe an ordinal as ``finite'' or ``infinite'' coincide:

\begin{lem}\ollabel{ordinfinitycharacter}
For any ordinal $\alpha$, the following are equivalent:
\begin{enumerate}
	\item\ollabel{ord:notinomega} $\alpha\notin \omega$, i.e.,
	$\alpha$ is not a natural number
	\item\ollabel{ord:omegaplus} $\omega \leq \alpha$ 
	\item\ollabel{ord:oneplus} $1 \ordplus  \alpha = \alpha$ 
	%\alpha \approx \alpha \ordplus 1$
	\item\ollabel{ord:plusone} $\alpha \approx \alpha\ordplus 1$,
	i.e., $\alpha$ and $\alpha\ordplus 1$ are equinumerous
	\item\ollabel{ord:infinite} $\alpha$ is Dedekind infinite	
	\end{enumerate}
\end{lem}
\noindent
So we have five provably equivalent ways to understand what it takes for an ordinal to be (in)finite.

\begin{proof}
\emph{\olref{ord:notinomega} $\Rightarrow$ \olref{ord:omegaplus}.} By
Trichotomy. 

\emph{\olref{ord:omegaplus} $\Rightarrow$ \olref{ord:oneplus}.} Fix
$\alpha \geq \omega$. By Transfinite Induction, there is some least
ordinal $\gamma$ (possibly $0$) such that there is a limit ordinal
$\beta$ with $\alpha = \beta \ordplus \gamma$. Now:
\[
	1 \ordplus \alpha =  
	1 \ordplus (\beta \ordplus \gamma) = 
	(1 \ordplus \beta) \ordplus \gamma =  
	\supstrict_{\delta < \beta} (1 \ordplus  \delta) \ordplus  \gamma = 
	\beta \ordplus  \gamma = 
	\alpha.
\]
\emph{\olref{ord:oneplus} $\Rightarrow$ \olref{ord:plusone}.} There is
clearly !!a{bijection} $f \colon (\alpha \disjointsum 1) \to (1
\disjointsum \alpha)$. If $1 \ordplus \alpha = \alpha$, there is an
isomorphism $g \colon (1 \disjointsum \alpha) \to \alpha$. Now
consider $\comp{f}{g}$.

\emph{\olref{ord:plusone} $\Rightarrow$ \olref{ord:infinite}.} If
$\alpha \approx \alpha \ordplus 1$, there is !!a{bijection} $f \colon
(\alpha \disjointsum 1) \to \alpha$. Define $g(\gamma) = f(\gamma, 0)$
for each $\gamma < \alpha$; this !!{injection} witnesses that $\alpha$
is Dedekind infinite, since $f(0,1) \in \alpha \setminus \ran{g}$. 

\emph{\olref{ord:infinite} $\Rightarrow$ \olref{ord:notinomega}.} This
is \olref[z][infinity-again]{naturalnumbersarentinfinite}.
\end{proof}

